\chapter{SABR et dynamique du smile}\label{ch:dv:sabr-and-smile-dynamics}

En 2002, quatre auteurs travaillant sur un desk de taux publient un court article aux conclusions
gênantes. Le modèle que de nombreux desks utilisent pour ajuster les smiles des swaptions, la
\omterm{def:dv:local-volatility:model}{volatilité locale}, prédit que lorsque le taux forward monte, le smile se déplace dans l'autre
sens, vers les strikes plus bas~; or le smile du marché suit le forward. Un modèle qui se trompe
de sens se trompe sur le delta, et les auteurs montrent que ses couvertures peuvent être pires
que l'absence totale de smile. Leur remplacement comporte quatre paramètres, une approximation
fermée de la volatilité implicite qu'un tableur sait évaluer, et la bonne dynamique. SABR devient
la norme de marché pour les options de taux et une référence courante pour les options de change.
Ce chapitre énonce le modèle et sa formule, explique le \omterm{def:dv:sabr-and-smile-dynamics:backbone}{backbone} qui gouverne sa dynamique, dérive
le delta qu'implique cette dynamique, mesure par simulation ce que vaut ce delta, et montre où la
formule cesse de tenir.

\section{Le modèle SABR}

\begin{definition}[Modèle SABR]\label{def:dv:sabr-and-smile-dynamics:sabr}
Le \emph{modèle SABR}\index{modèle SABR} (stochastic alpha, beta, rho) décrit un forward $F_t$,
sans dérive sous sa probabilité forward, et sa volatilité $\alpha_t$~:
\[
dF_t=\alpha_tF_t^{\beta}\,dW^1_t,\qquad d\alpha_t=\nu\,\alpha_t\,dW^2_t,\qquad
d\langle W^1,W^2\rangle_t=\rho\,dt,
\]
avec un exposant $\beta\in[0,1]$, une corrélation $\rho$ et une \omterm{def:dv:stochastic-volatility:sv}{volatilité de la volatilité} $\nu$~; la
valeur initiale $\alpha=\alpha_0$ constitue le quatrième paramètre.
\end{definition}

Il s'agit d'un \omterm{def:dv:stochastic-volatility:sv}{modèle à volatilité stochastique} au sens du chapitre~10, avec deux différences par
rapport à Heston~: la volatilité est lognormale et ne revient pas à la moyenne, ce qui convient à un
modèle calibré échéance par échéance, et le forward suit une diffusion à élasticité constante
$F^\beta$, qui interpole entre un modèle normal ($\beta=0$, le modèle de Bachelier du livre~2,
chapitre~13) et un modèle lognormal ($\beta=1$). Les desks de taux calibrent un SABR distinct pour
chaque échéance et chaque maturité du cube de volatilité~; les versions normales et décalées
qu'ils utilisent lorsque les taux approchent de zéro relèvent du livre~6, chapitre~5.

\section{La formule asymptotique}

SABR n'admet pas de prix en forme fermée. Ses auteurs ont dérivé, par perturbation singulière en
temps à l'échéance et en \omterm{def:dv:implied-volatility-and-its-surface:surface}{log-moneyness}, une approximation explicite de la volatilité implicite de Black.

\begin{definition}[Formule de Hagan]\label{def:dv:sabr-and-smile-dynamics:hagan}
La \emph{formule de Hagan}\index{formule de Hagan} de la volatilité de Black SABR au strike $K$ et à
l'échéance $T$ est
\[
\sigma_B(K)=\frac{\alpha}{(FK)^{\frac{1-\beta}2}\bigl(1+\frac{(1-\beta)^2}{24}\ell^2+
\frac{(1-\beta)^4}{1920}\ell^4\bigr)}\,\frac{z}{x(z)}\,\bigl(1+cT\bigr),
\]
\[
c=\frac{(1-\beta)^2\alpha^2}{24(FK)^{1-\beta}}+\frac{\rho\beta\nu\alpha}{4(FK)^{\frac{1-\beta}2}}
+\frac{2-3\rho^2}{24}\nu^2,
\]
avec $\ell=\ln(F/K)$, $z=\frac\nu\alpha(FK)^{\frac{1-\beta}2}\ell$ et
$x(z)=\ln\bigl((\sqrt{1-2\rho z+z^2}+z-\rho)/(1-\rho)\bigr)$.
\end{definition}

À la monnaie, $z/x(z)\to1$ et la formule se réduit à
$\sigma_{\mathrm{ATM}}=\alpha F^{\beta-1}\bigl(1+(\dots)T\bigr)$, une équation cubique en $\alpha$. Les desks
paramètrent donc SABR par la volatilité à la monnaie qu'ils observent, résolvent la cubique pour
$\alpha$, fixent $\beta$ par convention ou à partir du \omterm{def:dv:sabr-and-smile-dynamics:backbone}{backbone} (section suivante), puis ajustent $\rho$
et $\nu$ au smile. La volatilité normale (de Bachelier) que cote le marché des taux correspond au même prix, lu avec le modèle normal.

\begin{example}[Un forward à trois pour cent]\label{ex:dv:sabr-and-smile-dynamics:rates}
Une option à un an sur un taux forward de 3~\%, volatilité de Black à la monnaie de 20~\%, $\beta=0.5$, $\rho=-0.6$,
$\nu=0.5$~: la cubique donne $\alpha=0.0346$. Le smile de Black va de 28.8~\% pour un strike de 2~\% à 15.8~\% pour
4~\%~; les mêmes prix, lus en volatilités normales, vont de 70.8 à 55.0 points de base par an, avec 59.9 à la monnaie. Le
smile normal est beaucoup moins asymétrique~: une grande partie du skew de Black tient au facteur d'échelle propre au modèle lognormal.
\end{example}

\begin{example}[SABR sur un smile d'actions]\label{ex:dv:sabr-and-smile-dynamics:equity}
Avec $\beta=1$ (un \omterm{def:dv:sabr-and-smile-dynamics:backbone}{backbone} lognormal), SABR ajusté au smile à un an de la surface du chapitre~9
donne $\alpha=0.194$, $\rho=-0.68$, $\nu=0.72$, avec une erreur quadratique moyenne de 0.07 point de volatilité et au plus
0.15 point, au strike 70 (\cref{fig:dv:sabr-and-smile-dynamics:equity}). Une échéance, trois paramètres libres,
un bon ajustement~: SABR est autant un interpolateur de smile qu'un modèle.
\end{example}

\begin{omfigure}
\begin{tikzpicture}
\begin{axis}[width=10.5cm,height=4.8cm,
  xlabel={strike},ylabel={volatilité implicite à 1 an (\%)},
  tick label style={font=\small},label style={font=\small},
  xmin=68,xmax=132,ymin=14,ymax=30,grid=major,grid style={gray!25},
  legend columns=2,legend style={font=\footnotesize,at={(0.5,-0.3)},anchor=north,draw=none,fill=none,/tikz/every even column/.append style={column sep=8pt}},
  table/col sep=comma]
\addplot[omDef,only marks,mark=*,mark size=1.8pt] table[x=strike,y=market]{figdata/derivatives/11-sabr-and-smile-dynamics/equity_fit.csv};
\addplot[omThm,very thick] table[x=strike,y=sabr]{figdata/derivatives/11-sabr-and-smile-dynamics/equity_fit.csv};
\legend{marché (surface du chapitre~9),SABR avec $\beta=1$}
\end{axis}
\end{tikzpicture}

\omcaption{SABR à \omterm{def:dv:sabr-and-smile-dynamics:backbone}{backbone} lognormal ajusté à une échéance d'actions~: 13 strikes, erreurs d'au plus 0.15
point de volatilité. Données~: le tutoriel.}\label{fig:dv:sabr-and-smile-dynamics:equity}
\end{omfigure}

\section{Backbone et dynamique du smile}

\begin{definition}[Backbone]\label{def:dv:sabr-and-smile-dynamics:backbone}
Le \emph{backbone}\index{backbone} d'un modèle de smile est la courbe décrite par la volatilité implicite à la monnaie
lorsque le forward se déplace, toutes les autres variables d'état étant fixées. Dans SABR, elle vaut
$\sigma_{\mathrm{ATM}}(F)\approx\alpha F^{\beta-1}$~: plate pour $\beta=1$, décroissante comme $1/F$ pour $\beta=0$.
\end{definition}

Le \omterm{def:dv:sabr-and-smile-dynamics:backbone}{backbone} est la manière dont SABR encode la dynamique. Lorsque le forward bouge, le smile entier
le suit~: le point à la monnaie glisse le long du \omterm{def:dv:sabr-and-smile-dynamics:backbone}{backbone} et la forme du smile se déplace avec le
nouveau forward (\cref{fig:dv:sabr-and-smile-dynamics:backbone}). C'est le comportement que Hagan et ses coauteurs ont observé
sur les marchés de taux, et l'exact opposé de la \omterm{def:dv:local-volatility:model}{volatilité locale} (chapitre~9), dont le smile évolue à contre-courant du forward. Deux
paramètres produisent du skew~: $\beta<1$ donne un skew le long du \omterm{def:dv:sabr-and-smile-dynamics:backbone}{backbone}, $\rho<0$ un skew par rapport à lui. Ils
sont presque interchangeables pour ajuster un seul smile, et ne se distinguent que par la dynamique~: $\beta$ se
choisit d'après la manière dont la volatilité à la monnaie a évolué historiquement avec le forward, ou par convention.

\begin{omfigure}
\begin{tikzpicture}
\begin{groupplot}[group style={group size=2 by 1,horizontal sep=1.5cm},
  width=6.6cm,height=4.8cm,
  tick label style={font=\small},label style={font=\small},grid=major,grid style={gray!25},table/col sep=comma]
\nextgroupplot[xlabel={forward (\%)},ylabel={volatilité à la monnaie (\%)},xmin=2,xmax=4,ymin=14,ymax=31,
  title={\small backbones},
  legend columns=3,legend style={font=\footnotesize,at={(0.5,-0.32)},anchor=north,draw=none,fill=none}]
\addplot[omThm,very thick] table[x=fwd,y=b0]{figdata/derivatives/11-sabr-and-smile-dynamics/backbone.csv};
\addplot[omDef,thick] table[x=fwd,y=b05]{figdata/derivatives/11-sabr-and-smile-dynamics/backbone.csv};
\addplot[omProp,thick,dashed] table[x=fwd,y=b1]{figdata/derivatives/11-sabr-and-smile-dynamics/backbone.csv};
\legend{$\beta=0$,$\beta=0.5$,$\beta=1$}
\nextgroupplot[xlabel={strike (\%)},ylabel={volatilité de Black (\%)},xmin=1.5,xmax=5,ymin=14,ymax=36,
  title={\small smiles lorsque le forward bouge ($\beta=0.5$)},
  legend columns=3,legend style={font=\footnotesize,at={(0.5,-0.32)},anchor=north,draw=none,fill=none}]
\addplot[omThm,thick] table[x=k,y=f025]{figdata/derivatives/11-sabr-and-smile-dynamics/smiles.csv};
\addplot[omDef,very thick] table[x=k,y=f03]{figdata/derivatives/11-sabr-and-smile-dynamics/smiles.csv};
\addplot[omProp,thick,dashed] table[x=k,y=f035]{figdata/derivatives/11-sabr-and-smile-dynamics/smiles.csv};
\legend{$F=2.5\%$,$F=3\%$,$F=3.5\%$}
\end{groupplot}
\end{tikzpicture}

\omcaption{Dynamique du SABR. À gauche~: la volatilité à la monnaie en fonction du forward, pour trois valeurs de
$\beta$ avec $\alpha$ fixé de sorte à obtenir 20~\% pour un forward de 3~\%. À droite~: le smile pour trois forwards ($\beta=0.5$,
$\rho=-0.6$, $\nu=0.5$)~: il se déplace avec le forward, son minimum se déplaçant vers la droite lorsque le forward monte.
Données~: le code du chapitre.}\label{fig:dv:sabr-and-smile-dynamics:backbone}
\end{omfigure}

\section{Couverture sous un modèle de smile}

Le delta d'un modèle de smile est la variation de la valeur de l'option lorsque le sous-jacent varie, avec la vision
qu'a le modèle de l'évolution des autres variables d'état dans l'intervalle. Trois réponses sont utilisées.

\begin{itemize}
\item Le \emph{delta de Black} gèle la volatilité implicite de l'option~: $\Delta_{\mathrm{BS}}$.
\item Le \emph{delta de Hagan} déplace la volatilité le long du smile lorsque $F$ varie, avec $\alpha$ fixé~:
$\Delta_{\mathrm{BS}}+\mathcal V\,\partial\sigma_B/\partial F$.
\item Le troisième tient compte du fait que $\alpha$ lui-même tend à varier lorsque $F$ varie, puisque les deux
mouvements browniens sont corrélés.
\end{itemize}

\begin{definition}[Delta de variance minimale]\label{def:dv:sabr-and-smile-dynamics:mvdelta}
Le \emph{delta de variance minimale}\index{delta de variance minimale} d'une option dans un \omterm{def:dv:stochastic-volatility:sv}{modèle à volatilité stochastique}
est la position sur le sous-jacent qui minimise la variance du P\&L instantané de la position couverte.
Dans le SABR (delta de Bartlett), il vaut
\[
\Delta_{\mathrm{MV}}=\Delta_{\mathrm{BS}}+\mathcal V\Bigl(\frac{\partial\sigma_B}{\partial F}
+\frac{\partial\sigma_B}{\partial\alpha}\,\frac{\rho\nu}{F^{\beta}}\Bigr).
\]
\end{definition}

\begin{proposition}[Pourquoi ce delta]\label{prop:dv:sabr-and-smile-dynamics:mv}
La valeur de l'option $V(F,\alpha)$ varie de $V_F\,dF+V_\alpha\,d\alpha$ au premier ordre. La régression de
$d\alpha$ sur $dF$ a pour pente $\rho\nu/F^\beta$, de sorte que la position $\Delta=V_F+V_\alpha\rho\nu/F^\beta$ élimine
la partie de $d\alpha$ corrélée à $dF$ et ne laisse subsister que la partie non corrélée, qu'aucune position sur
le sous-jacent ne peut couvrir.
\end{proposition}
\begin{proof}
$d\alpha=\nu\alpha\,dW^2=\nu\alpha(\rho\,dW^1+\sqrt{1-\rho^2}\,dW^\perp)$ et $dF=\alpha F^\beta dW^1$, donc
$d\alpha=\frac{\rho\nu}{F^\beta}dF+\nu\alpha\sqrt{1-\rho^2}\,dW^\perp$. En substituant, le P\&L de
$V-\Delta F$ vaut $(V_F+V_\alpha\rho\nu/F^\beta-\Delta)\,dF+V_\alpha\nu\alpha\sqrt{1-\rho^2}\,dW^\perp$~; sa
variance est minimale lorsque le premier crochet s'annule. Avec $V=C_{\mathrm{BS}}(F,\sigma_B(F,\alpha))$,
$V_F=\Delta_{\mathrm{BS}}+\mathcal V\partial_F\sigma_B$ et $V_\alpha=\mathcal V\partial_\alpha\sigma_B$.
\end{proof}

Sur l'exemple du chapitre, le call à un an à la monnaie a un delta de Black de 0.540, un delta de Hagan de 0.580
et un \omterm{def:dv:sabr-and-smile-dynamics:mvdelta}{delta de variance minimale} de 0.461 (\cref{fig:dv:sabr-and-smile-dynamics:deltas}). Avec une corrélation négative,
une hausse du forward tend à entraîner une baisse de $\alpha$, donc de la valeur du call~: la bonne couverture est plus petite
que celle de Black. Le delta de Hagan évolue dans l'autre sens parce qu'il ignore la corrélation et ne voit que le \omterm{def:dv:sabr-and-smile-dynamics:backbone}{backbone}.

\begin{omfigure}
\begin{tikzpicture}
\begin{axis}[width=10.5cm,height=4.8cm,
  xlabel={strike (\%)},ylabel={delta d'un call à 1 an},
  tick label style={font=\small},label style={font=\small},
  xmin=2,xmax=4.5,ymin=0,ymax=1,grid=major,grid style={gray!25},
  legend columns=3,legend style={font=\footnotesize,at={(0.5,-0.3)},anchor=north,draw=none,fill=none,/tikz/every even column/.append style={column sep=8pt}},
  table/col sep=comma]
\addplot[gray,thick,dashed] table[x=k,y=black]{figdata/derivatives/11-sabr-and-smile-dynamics/deltas.csv};
\addplot[omProp,thick] table[x=k,y=hagan]{figdata/derivatives/11-sabr-and-smile-dynamics/deltas.csv};
\addplot[omThm,very thick] table[x=k,y=bartlett]{figdata/derivatives/11-sabr-and-smile-dynamics/deltas.csv};
\legend{Black,Hagan,variance minimale}
\end{axis}
\end{tikzpicture}

\omcaption{Trois deltas de calls à un an sur un forward de 3~\% en fonction du strike ($\beta=0.5$, $\rho=-0.6$, $\nu=0.5$).
Le \omterm{def:dv:sabr-and-smile-dynamics:mvdelta}{delta de variance minimale} est le plus petit, car une hausse du forward est censée s'accompagner d'une baisse de
volatilité. Données~: le tutoriel.}\label{fig:dv:sabr-and-smile-dynamics:deltas}
\end{omfigure}

Quelle est l'importance de cet effet~? Le tutoriel simule une journée de dynamique SABR sur 200\,000 trajectoires et réévalue le
call avec la formule au nouveau forward et au nouvel $\alpha$. L'écart-type du P\&L couvert sur une journée vaut
$6.67\times10^{-5}$ avec le delta de Black, $7.46\times10^{-5}$ avec celui de Hagan et $5.99\times10^{-5}$ avec le
\omterm{def:dv:sabr-and-smile-dynamics:mvdelta}{delta de variance minimale}~: une variance inférieure de 19~\% à celle de Black, et supérieure de 25~\% avec celui de Hagan. Il reste la part du
mouvement de la volatilité indépendante du forward, que l'on ne couvre qu'avec le \omterm{def:dv:greeks-and-the-hedging-pnl:greeks}{vega} (une autre option).

\section{Limites~: ailes et échéances longues}

La \omterm{def:dv:sabr-and-smile-dynamics:hagan}{formule de Hagan} est un développement en $T$ et en $\ln(F/K)$. Elle est précise pour les échéances et les strikes du
marché liquide des swaptions, et elle échoue de deux façons connues. Aux échéances longues avec une forte \omterm{def:dv:stochastic-volatility:sv}{volatilité de la
volatilité}, elle produit des prix d'options qui ne sont pas convexes pour les bas strikes~: une densité négative
(\cref{fig:dv:sabr-and-smile-dynamics:density}). Son terme d'ordre zéro s'est par ailleurs révélé légèrement faux
(correction d'Obłój), ce qui compte dans les ailes. Les correctifs utilisés en pratique consistent à résoudre une équation
forward unidimensionnelle effective pour la densité, avec la même précision asymptotique («~SABR sans arbitrage~»), ou
à évaluer les ailes avec une autre extrapolation, ou à décaler le forward (le livre~6, chapitre~5).

\begin{omfigure}
\begin{tikzpicture}
\begin{axis}[width=10.5cm,height=4.8cm,
  xlabel={strike (\%)},ylabel={densité issue des prix de Hagan},
  tick label style={font=\small},label style={font=\small},
  xmin=0,xmax=3,ymin=-130,ymax=60,grid=major,grid style={gray!25},
  table/col sep=comma]
\addplot[omThm,very thick] table[x=k,y=density]{figdata/derivatives/11-sabr-and-smile-dynamics/density.csv};
\draw[gray,thick] (axis cs:0,0) -- (axis cs:3,0);
\node[font=\footnotesize,anchor=west,align=left] at (axis cs:1.7,-80) {densité négative~:\\un arbitrage butterfly};
\end{axis}
\end{tikzpicture}

\omcaption{La densité de Breeden--Litzenberger issue de la \omterm{def:dv:sabr-and-smile-dynamics:hagan}{formule de Hagan} pour une option à dix ans sur un forward de 3~\%
($\beta=0.5$, $\rho=-0.3$, $\nu=0.6$, 20~\% à la monnaie). En dessous d'environ 1.6~\% elle est négative~: la formule donne
à ces butterflies des prix négatifs. Données~: le tutoriel.}\label{fig:dv:sabr-and-smile-dynamics:density}
\end{omfigure}

\section{Tutoriel~: SABR de la formule à la couverture}\label{tut:dv:sabr-and-smile-dynamics:tut}

\begin{tutorial}
\textbf{Objectif.} Implémenter la \omterm{def:dv:sabr-and-smile-dynamics:hagan}{formule de Hagan}, calibrer le SABR avec $\alpha$ lié à la volatilité à la monnaie, comparer
trois deltas et mesurer leurs erreurs de couverture par simulation. \textbf{État final~:} les quatre figures du chapitre
et les nombres du problème du week-end.
\begin{tutsteps}
\item \textbf{La formule}, avec la limite en la monnaie de $z/x(z)$ traitée explicitement~:
\omcode{code/firm/sabr/firm_sabr.py}{23}{37}{Formule de Hagan pour la volatilité de Black.}\label{lst:dv:sabr-and-smile-dynamics:hagan}
\item \textbf{Les trois deltas}, la correction de corrélation provenant de $\partial\sigma_B/\partial\alpha$~:
\omcode{code/firm/sabr/firm_sabr.py}{76}{89}{Deltas de Black, de Hagan et de variance minimale.}\label{lst:dv:sabr-and-smile-dynamics:deltas}
\item \textbf{Exécuter} \texttt{dv\_sabr.hedge\_experiment()}, \texttt{equity\_fit()}, \texttt{wing\_density()} et
\texttt{fig\_sabr.py}.
\end{tutsteps}
\textbf{Pour aller plus loin.} Poser $\rho=+0.3$ et comparer les erreurs de couverture des trois deltas (le \omterm{def:dv:sabr-and-smile-dynamics:mvdelta}{delta de variance
minimale} est alors presque égal à celui de Black)~; puis ajuster le smile des actions avec $\beta=0.5$ et comparer $\rho$ et $\nu$ avec le
calibrage à $\beta=1$.
\end{tutorial}

\section{Construction~: le smile SABR}\label{bld:dv:sabr-and-smile-dynamics:sabr}

\begin{build}
\bfield{Rôle} Le smile de la firme miniature pour les options de taux et de change, et le paramétrage que le
livre~6 étend au cube de volatilité.
\bfield{Interface} \texttt{hagan\_lognormal(f, k, t, alpha, beta, rho, nu)}~; \texttt{alpha\_from\_atm(f, t,
sigma\_atm, beta, rho, nu)}~; \texttt{calibrate(f, t, strikes, vols, beta, atm\_vol) -> ((alpha, rho, nu), rmse)}~;
\texttt{normal\_vol}~; \texttt{deltas(\ldots) -> black, hagan, bartlett}~; \texttt{one\_day\_hedge\_errors}.
\bfield{Règles} $\beta$ est une donnée d'entrée, et non un paramètre calibré~; $\alpha$ est toujours résolu à partir de la volatilité à la monnaie
(la plus petite racine positive du cubique)~; il faut présenter la densité des ailes, et ne pas utiliser la formule là où elle est
négative.
\bfield{Tests d'acceptation} \texttt{code/firm/sabr/tests/}~: pour $\nu\to0$ et $\beta=1$, le smile est plat à $\alpha$~; la
volatilité à la monnaie est reproduite exactement~; le calibrage retrouve des paramètres connus~; le \omterm{def:dv:sabr-and-smile-dynamics:mvdelta}{delta de variance minimale} présente le
plus faible erreur de couverture simulée pour les corrélations négatives.
\bfield{Pour aller plus loin} La densité sans arbitrage de Hagan et de ses coauteurs (2014), résolue sur une grille, et le terme
principal corrigé d'Obłój.
\end{build}

\begin{omsources}
\item P.~S.~Hagan, D.~Kumar, A.~S.~Lesniewski et D.~E.~Woodward, «~Managing smile risk~», \emph{Wilmott} (2002).
\item B.~Bartlett, «~Hedging under SABR model~», \emph{Wilmott} (juillet/août 2006) 2--4.
\item J.~Obłój, «~Fine-tune your smile: correction to Hagan et al.~», \emph{Wilmott} (mai 2008) 102--109.
\item P.~S.~Hagan, D.~Kumar, A.~S.~Lesniewski et D.~E.~Woodward, «~Arbitrage-free SABR~», \emph{Wilmott} 69 (2014)
60--75.
\end{omsources}

\section{Exercices}

\begin{exercise}[$\star$]\label{exo:dv:sabr-and-smile-dynamics:1}
Avec $\beta=1$ et $\nu=0$, quel est le smile SABR~? Et avec $\beta=0$ et $\nu=0$~?
\end{exercise}

\begin{exercise}[$\star$]\label{exo:dv:sabr-and-smile-dynamics:2}
En utilisant le terme principal $\sigma_{\mathrm{ATM}}\approx\alpha F^{\beta-1}$ avec $\alpha$ fixé pour 20~\% à un forward de 3~\%, donner
la volatilité à la monnaie à un forward de 2.5~\% pour $\beta=0$, 0.5 et 1.
\end{exercise}

\begin{exercise}[$\star$]\label{exo:dv:sabr-and-smile-dynamics:3}
Pourquoi SABR n'a-t-il pas besoin de retour à la moyenne lorsqu'il est calibré échéance par échéance~?
\end{exercise}

\begin{exercise}[$\star\star$]\label{exo:dv:sabr-and-smile-dynamics:4}
Dériver le terme de correction du \omterm{def:dv:sabr-and-smile-dynamics:mvdelta}{delta de variance minimale} à partir de la régression de $d\alpha$ sur $dF$.
\end{exercise}

\begin{exercise}[$\star\star$]\label{exo:dv:sabr-and-smile-dynamics:5}
$\beta$ et $\rho$ produisent tous deux du skew. Comment choisiriez-vous $\beta$ pour une paire de devises, et que
coûte le mauvais choix~?
\end{exercise}

\begin{exercise}[$\star\star$]\label{exo:dv:sabr-and-smile-dynamics:6}
Donner la volatilité normale des strikes à 2~\%, 3~\% et 4~\% dans l'\cref{ex:dv:sabr-and-smile-dynamics:rates}, et expliquer pourquoi
le smile normal est plus plat que le smile de Black.
\end{exercise}

\begin{exercise}[$\star\star\star$]\label{exo:dv:sabr-and-smile-dynamics:7}
\emph{Codage.} Reprendre l'expérience de couverture à un jour avec $\rho=+0.3$, $\nu=0.4$~: quel delta est le meilleur, et de combien~?
\end{exercise}

\begin{exercise}[$\star\star\star$]\label{exo:dv:sabr-and-smile-dynamics:8}
\emph{Repérer l'erreur.} «~Notre calibrage SABR est parfait à toutes les échéances, donc son prix d'exercice à 1~\% à dix ans est fiable.~»
\end{exercise}

\section{Problème~: quel delta}

\begin{problem}[{Problème du week-end --- le delta qu'un modèle de smile doit utiliser}]\label{pb:dv:sabr-and-smile-dynamics:1}
Un desk de taux est long d'options d'achat à un an à la monnaie sur un forward à 3~\% (swaptions payeuses en unités normalisées), évaluées
avec SABR~: $\beta=0.5$, $\rho=-0.6$, $\nu=0.5$, volatilité de Black à la monnaie de 20~\%, actualisation nulle. Il doit choisir le
delta avec lequel il se couvre.

\medskip\noindent\textbf{Partie I --- Le modèle.}
\begin{enumerate}
\item Résoudre pour $\alpha$.
\item Donner la volatilité de Black aux strikes à 2~\%, 3~\% et 4~\%.
\item Donner la volatilité normale à la monnaie en points de base.
\item Que prédit le \omterm{def:dv:sabr-and-smile-dynamics:backbone}{backbone} pour la volatilité à la monnaie si le forward monte à 3.5~\%~?
\item Comment le smile se déplace-t-il lorsque le forward monte~?
\end{enumerate}

\medskip\noindent\textbf{Partie II --- Trois deltas.}
\begin{enumerate}[resume]
\item Donner les deltas de Black, de Hagan et de variance minimale de l'option d'achat à la monnaie.
\item Pourquoi le delta de Hagan est-il plus grand que celui de Black~?
\item Pourquoi le \omterm{def:dv:sabr-and-smile-dynamics:mvdelta}{delta de variance minimale} est-il plus petit que les deux~?
\item Donner la pente de régression de $d\alpha$ sur $dF$.
\item Quel risque subsiste après la couverture de variance minimale, et comment le couvrir~?
\end{enumerate}

\medskip\noindent\textbf{Partie III --- L'expérience.}
\begin{enumerate}[resume]
\item Donner les écarts-types du \omterm{def:dv:greeks-and-the-hedging-pnl:hpnl}{P\&L de couverture} à un jour avec chaque delta.
\item Donner la variance de chacun par rapport à celle de Black.
\item Pourquoi le delta de Hagan fait-il moins bien que celui de Black ici~?
\item Quel serait le résultat avec $\rho=0$~?
\item Si les erreurs de couverture quotidiennes sont indépendantes, de combien une variance quotidienne inférieure de 19~\% réduit
l'écart-type du \omterm{def:dv:greeks-and-the-hedging-pnl:hpnl}{P\&L de couverture} annuel~?
\end{enumerate}

\medskip\noindent\textbf{Partie IV --- Jugement.}
\begin{enumerate}[resume]
\item Quel delta faudrait-il présenter au risque, et lequel faudrait-il trader~?
\item Comment le choix de $\beta$ modifie-t-il les deltas~?
\item Que se passe-t-il si la dynamique du smile de marché diffère de celle de SABR~?
\item Énoncer le \emph{résultat nommé}~: la réduction de la variance de l'erreur de couverture à un jour
due au \omterm{def:dv:sabr-and-smile-dynamics:mvdelta}{delta de variance minimale} par rapport au delta de Black.
\item En une phrase~: de quoi dépend le delta d'un modèle de smile, qui ne figure pas dans ses prix~?
\end{enumerate}
\end{problem}

\section{Questions d'entretien}

\begin{interviewq}[$\star$ \iqroles{trader, chercheur}]\label{iq:dv:sabr-and-smile-dynamics:1}
Écrivez le \omterm{def:dv:sabr-and-smile-dynamics:sabr}{modèle SABR}. Que fait chaque paramètre~?
\end{interviewq}

\begin{interviewq}[$\star$ \iqroles{trader}]\label{iq:dv:sabr-and-smile-dynamics:2}
Qu'est-ce que le \omterm{def:dv:sabr-and-smile-dynamics:backbone}{backbone}, et comment estimeriez-vous $\beta$ à partir des données de marché~?
\end{interviewq}

\begin{interviewq}[$\star\star$ \iqroles{chercheur}]\label{iq:dv:sabr-and-smile-dynamics:3}
Dérivez le \omterm{def:dv:sabr-and-smile-dynamics:mvdelta}{delta de variance minimale} dans un \omterm{def:dv:stochastic-volatility:sv}{modèle à volatilité stochastique}.
\end{interviewq}

\begin{interviewq}[$\star\star$ \iqroles{chercheur, risque}]\label{iq:dv:sabr-and-smile-dynamics:4}
Pourquoi la \omterm{def:dv:local-volatility:model}{volatilité locale} a-t-elle donné de mauvaises couvertures sur le marché des taux, selon les auteurs de SABR~?
\end{interviewq}

\begin{interviewq}[$\star\star$ \iqroles{développeur}]\label{iq:dv:sabr-and-smile-dynamics:5}
Votre implémentation de SABR renvoie des prix de butterfly négatifs pour les strikes bas à longue échéance. Pourquoi, et que faites-vous~?
\end{interviewq}

\begin{interviewq}[$\star\star\star$ \iqroles{trader, chercheur}]\label{iq:dv:sabr-and-smile-dynamics:6}
$\beta$ et $\rho$ ajustent le même smile aussi bien l'un que l'autre. Quelles transactions permettent de les distinguer, et comment
couvririez-vous un livre dont ils sont incertains~?
\end{interviewq}
